\label{chap:p3-83-45-completitud_finita}

\section{Extension conservative de la complétude finie par des multiplicités projectives}

Une tour HMT concrète est formée d'ensembles finis, de relations de
transition, d'applications de restriction et de lecteurs. Tout cela peut
être codé dans ZF ou ZFC. La collection de toutes les tours constitue une
classe définissable ; il n'est pas nécessaire de postuler que HMT forme déjà une théorie
autonome des classes comparable à NBG.

\begin{proposition}[Codage conservatif]
Chaque objet généalogique fini ou dénombrable utilisé dans ce cahier peut
être représenté par des ensembles et des relations de ZFC. Par conséquent, les
formalisations construites ici ne démontrent pas une contradiction de ZFC.
\end{proposition}

\begin{proof}
Chaque niveau \(X_m\) est un ensemble, chaque restriction \(p_{n,m}\) est un
sous-ensemble de \(X_n\times X_m\), et une tour dénombrable est une fonction de
domaine \(\mathbb N\) qui code ces données. Les types, lecteurs et mémoires
s'ajoutent comme relations ou fonctions. La limite inverse est un sous-ensemble
du produit \(\prod_mX_m\), que ZFC construit par Remplacement, Union,
Ensemble des parties et Séparation.
\end{proof}

Ce que HMT discute n'est pas la légitimité de ce codage, mais la pratique
qui consiste à oublier les applications et à ne conserver qu'une ombre. Si l'on applique, par exemple,
le lecteur
\[
  L:\OmegaInf\longrightarrow\mathbb R
\]
ou le profil cardinal
\[
  \mathbb X\longmapsto(|X_m|)_{m\ge0},
\]
des objets non isomorphes peuvent être confondus. La projection est valide dans son
codomaine ; ce qui serait invalide serait de lui attribuer une fidélité sans preuve.

\begin{remark}[Conditions de validité]
Il n'est pas affirmé que HMT satisfasse en interne tous les schémas de ZF, IZF ou
NBG. Séparation, Remplacement, Ensemble des parties complet et une hiérarchie transfinie générale
n'ont pas été vérifiés conjointement dans une axiomatique HMT autonome.
\end{remark}

\section{Catégorie d'objets généalogiques}

\begin{definition}
Un objet généalogique est
\[
  \mathbb X=
  \bigl((X_m)_{m\ge0},(p_{n,m})_{n\ge m},
  \mathcal T,\mathcal O,\mathcal M\bigr),
\]
où \(\mathcal T\) enregistre les types, \(\mathcal O\) les observables et
\(\mathcal M\) les mémoires. Un morphisme
\(f:\mathbb X\to\mathbb Y\) est une famille \(f_m:X_m\to Y_m\) qui satisfait
\[
  p^Y_{n,m}f_n=f_mp^X_{n,m}
\]
et préserve les données indiquées.
\end{definition}

La catégorie ainsi obtenue est une présentation concrète d'objets
profinis enrichis. Le foncteur des sections globales
\[
  \Gamma(\mathbb X)=\varprojlim_mX_m
\]
retrouve les histoires complètes. La petite catégorie de chemins de
\(\TPK\) admet en outre l'enveloppe de préfaisceaux
\[
  \widehat{\mathcal P_{\TPK}}
  =
  \mathbf{Set}^{\mathcal P_{\TPK}^{op}}.
\]
Cette construction fournit une logique contextuelle et des faisceaux de données, mais
n'autorise pas à déclarer la booléanité, le choix ou l'existence de suffisamment de points sans étudier
le topos concret.

\section{Fondation des préfixes et coinduction des sections}

Le graphe des préfixes est gradué par la profondeur :
\[
  \operatorname{rk}(x_m)=m.
\]
Appliquer le parent diminue strictement le rang. Il n'existe pas de cycles
parentaux
\[
  x_m\to x_{m-1}\to\cdots\to x_m
\]
compatibles avec cette graduation.

\begin{proposition}
La bonne fondation du graphe des préfixes et l'existence coinductive d'une
section infinie sont compatibles.
\end{proposition}

\begin{proof}
La bonne fondation concerne la direction de descente par les parents. La
coinduction concerne une famille \((x_m)\) compatible avec toutes les
descentes. Aucun \(x_m\) n'appartient à lui-même ni ne réapparaît comme le même
nœud ; la section est une fonction sur \(\mathbb N\), et non un cycle
d'appartenance.
\end{proof}

Cette distinction éclaire le retour \(9\to1\). La phase peut se répéter, mais
la coordonnée de mémoire augmente. Une boucle dans la projection de phase n'est pas un
cycle dans le graphe gradué des états.

\section{Fibration des ordinaux par des histoires généalogiques compatibles}

Von Neumann représente chaque ordinal comme l'ensemble de ses prédécesseurs
\citep{vonneumann1923}. HMT peut relever cette représentation sans la modifier.

\begin{definition}[Relèvement généalogique d'un ordinal fini]
Soit \(\boldsymbol x=(x_m)\) une section. On définit
\[
  \widehat n_{\boldsymbol x}
  =
  \{(k,x_k):k<n\}.
\]
\end{definition}

\begin{proposition}
La projection sur la première coordonnée satisfait
\[
  \pi(\widehat n_{\boldsymbol x})=n,
\]
et la suite obéit à
\[
  \widehat{n+1}_{\boldsymbol x}
  =
  \widehat n_{\boldsymbol x}\cup\{(n,x_n)\}.
\]
\end{proposition}

\begin{proof}
Les deux identités s'obtiennent directement de la définition. L'information
ajoutée ne modifie pas l'ordre sous-jacent ; elle distingue les routes qui réalisent le
même ordinal.
\end{proof}

Pour \(\omega\), on écrit
\[
  \widehat\omega_{\boldsymbol x}
  =
  \{(n,x_n):n\in\mathbb N\}.
\]
Le support de chaque histoire reste dénombrable, bien que l'espace de
toutes les histoires puisse avoir la puissance du continu lorsque la ramification
persiste.

\begin{remark}[Extension et domaine ouvert]
L'extension
\[
S_{\alpha+1}=\operatorname{Ext}_\alpha(S_\alpha),
\qquad
S_\lambda=\varprojlim_{\beta<\lambda}S_\beta
\]
est un programme transfini bien typé, et non un résultat déjà démontré pour
tout ordinal. Il faut vérifier la cohérence des limites, la taille des
niveaux et l'existence d'extensions.
\end{remark}

\section{Foncteur de décatégorification et cardinalité des sections publiées}

Deux tours peuvent satisfaire
\[
  |X_m|=|Y_m|\qquad\text{pour tout }m
\]
et ne pas être isomorphes : il suffit que leurs applications de restriction aient des schémas de
ramification distincts. Le profil cardinal est donc un invariant
utile mais non fidèle.

\begin{remark}[Interprétation structurelle]
Compter combien de cases existent ne dit pas quelle case en prolonge quelle autre, quelle
orientation s'inverse, quelle retenue est transportée ou quel état revient. Le
cardinal mesure l'inventaire ; la généalogie mesure l'organisation.
\end{remark}

\begin{remark}[Séparation des résultats et portée]
Le codage ZFC, la section comme limite et la distinction entre égalité
observationnelle et identité d'état sont des
\emph{résultats établis dans les sections précédentes}. La catégorie unifiée, le topos de
préfaisceaux et l'ordinal fibré sont des
\emph{constructions explicites du présent développement}. Il n'est pas proclamé de théorie autonome complète
des classes ou des ordinaux transfinis HMT.
\end{remark}


\section{Limite profinie du système de troncatures généalogiques}

Soit
\[
  X=\varprojlim_mX_m
\]
une limite inverse d'ensembles finis discrets avec des applications surjectives.
Nous munissons \(X\) de la topologie initiale des projections
\(\pi_m:X\to X_m\). Les cylindres
\[
  [A]_m=\pi_m^{-1}(A),
  \qquad A\subseteq X_m,
\]
sont ouverts-fermés et forment une base.

\begin{theorem}[Ensemble des parties cylindriques]
\label{p3-83-c45-s2:thm:clopen-power}
\[
  \Clop(X)
  \cong
  \varinjlim_m\bigl(\mathcal P(X_m),p_{n,m}^{-1}\bigr).
\]
\end{theorem}

\begin{proof}
Tout cylindre est ouvert-fermé. Réciproquement, si \(C\subset X\) est ouvert-fermé, pour
chaque \(x\in C\) il existe un cylindre contenu dans \(C\). La compacité de \(C\)
donne un sous-recouvrement fini. En relevant ces cylindres à une profondeur
commune \(m\), on obtient \(C=\pi_m^{-1}(A)\) pour un certain
\(A\subseteq X_m\). Deux représentants à des profondeurs distinctes définissent le
même ouvert-fermé exactement lorsqu'ils coïncident après un relèvement commun.
\end{proof}

\begin{theorem}[Puissance sous forme close]
\label{p3-83-c45-s2:thm:closed-power}
\[
  \Closed(X)
  \cong
  \varprojlim_m\bigl(\mathcal P(X_m),p_{n,m*}\bigr),
\]
où \(p_*(A)=p(A)\).
\end{theorem}

\begin{proof}
Si \(F\subset X\) est fermé, il est compact. Ses projections
\(F_m=\pi_m(F)\) satisfont \(p_{n,m}(F_n)=F_m\). Réciproquement, une famille
compatible \((F_m)\) définit la limite inverse
\[
  F=\{x\in X:\pi_m(x)\in F_m\text{ pour tout }m\}.
\]
C'est une intersection de fermés et, par la compatibilité surjective des
restrictions \(F_n\to F_m\), ses projections sont exactement les \(F_m\).
Les deux constructions sont inverses.
\end{proof}

\section{Obstruction à la reconstruction de l'ensemble des parties généalogiques à partir de réductions finies}

Dans un espace de Cantor,
\[
 |\Clop(X)|=\aleph_0,
 \qquad
 |\Closed(X)|=2^{\aleph_0},
 \qquad
 |\mathcal P(X)|=2^{2^{\aleph_0}}.
\]
Les trois objets ne peuvent pas être identifiés.

Cette séparation ne présuppose pas l'hypothèse du continu. En écrivant
\(\mathfrak c=2^{\aleph_0}\),
\[
 |\Clop(X)|=\aleph_0,\qquad
 |\Closed(X)|=\mathfrak c,\qquad
 |\mathcal P(X)|=2^{\mathfrak c}.
\]
\(\mathsf{CH}\) identifie seulement \(\mathfrak c\) à \(\aleph_1\) ; elle ne
confond aucune des trois couches. L'indépendance de Gödel--Cohen change quel
cardinal peut occuper \(\mathfrak c\), et non l'inégalité de Cantor
\(\mathfrak c<2^{\mathfrak c}\).

\begin{proposition}[Falsificateur dense]
Si \(D\subsetneq X\) est dense, alors
\[
  \pi_m(D)=X_m
  \qquad\text{pour tout }m,
\]
de sorte que \(D\) et \(X\) possèdent les mêmes ombres finies.
\end{proposition}

\begin{proof}
Chaque fibre cylindrique non vide \(\pi_m^{-1}(a)\) est ouverte. Par densité,
elle rencontre \(D\). Par conséquent, \(a\in\pi_m(D)\) pour tout \(a\in X_m\).
\end{proof}

Les ombres finies reconstruisent l'adhérence d'un sous-ensemble, et non le
sous-ensemble arbitraire. C'est le point exact où une identification
\[
  \varprojlim_m\mathcal P(X_m)
  =
  \mathcal P(\varprojlim_mX_m)
\]
serait fausse.

\begin{figure}[H]
\centering
\begin{tikzpicture}[>=Latex,node distance=10mm and 17mm,
  box/.style={rounded corners,draw=hmtblue,fill=hmtlightblue,
    minimum width=38mm,minimum height=15mm,align=center},
  full/.style={rounded corners,draw=hmtviolet,fill=white,
    minimum width=40mm,minimum height=15mm,align=center}]
  \node[box] (fin) {\(\mathcal P(X_m)\)\\\scriptsize prédicats finis};
  \node[box,right=of fin] (clop) {\(\Clop(X)\)\\\scriptsize limite directe};
  \node[box,below=of clop] (closed) {\(\Closed(X)\)\\\scriptsize limite inverse};
  \node[full,right=of clop] (power) {\(\mathcal P(X)\)\\\scriptsize ensemble complet des parties};
  \draw[->,thick,hmtblue] (fin)--node[above]{\(p^{-1}\)}(clop);
  \draw[->,thick,hmtblue] (fin) to[bend right=24]
    node[below left]{\(p_*\)} (closed);
  \draw[->,thick,hmtgreen] (clop)--node[above]{\(\subset\)}(power);
  \draw[->,thick,hmtgreen] (closed)--node[below]{\(\subset\)}(power);
  \draw[dashed,very thick,hmtred] (closed.south east)
    .. controls +(8mm,-8mm) and +(-8mm,-8mm) ..
    node[midway,below=2mm]{\(\ne\)} (power.south west);
  \node[below=18mm of closed,text width=.82\textwidth,align=center] {
    \small Un sous-ensemble dense propre et \(X\) ont les mêmes projections finies :
    l'histoire finie retrouve l'adhérence, et non l'ensemble complet des parties.};
\end{tikzpicture}
\caption{Trois notions d'ensemble des parties qu'il ne faut pas aplatir.}
\label{p3-83-c45-s2:fig:three-powers}
\end{figure}

\section{\(4\) exigences de choix dans le système projectif}

Soit \(p:X\twoheadrightarrow Y\). Il convient de séparer :

\begin{enumerate}
\item une \emph{section ensembliste} \(s:Y\to X\) avec \(ps=\id_Y\) ;
\item une \emph{famille naturelle} de sections qui commute avec les
restrictions d'une tour ;
\item une \emph{section algébrique} qui préserve en outre les opérations ;
\item une \emph{section effective} calculée par un algorithme total.
\end{enumerate}

Toute section naturelle, algébrique ou effective est, en oubliant la structure, une
section ensembliste. Les trois propriétés supplémentaires sont, en général,
incomparables : un algorithme peut choisir des représentants sans être un homomorphisme ;
un homomorphisme peut ne pas être calculable dans la présentation choisie ; une
famille naturelle peut ne pas conserver l'opération algébrique.

\begin{remark}[Conditions de validité]
L'axiome du choix garantit une sélection ensembliste dans son
domaine. Il ne garantit ni naturalité, ni algébricité, ni effectivité, ni coût
uniforme. L'indépendance de \(\mathsf{AC}\) par rapport à \(\mathsf{ZF}\)
ne transforme pas non plus ces quatre exigences en une seule : même dans un modèle
avec choix, la section structurelle nécessite sa propre preuve.
\end{remark}

\Needspace{36\baselineskip}
\section{Obstruction de Pell à la périodicité uniforme du système de préfixes}

Le générateur actif construit
\[
  K_\infty\cong C_4\times\mathbb Z_3
\]
et des projections finies \(p_m:K_\infty\to K_m\). Il est possible de choisir un
représentant de chaque classe, mais pas une section homomorphe compatible avec le
groupe.

\begin{proposition}
Il n'existe pas d'homomorphisme \(s_m:K_m\to K_\infty\) avec
\(p_ms_m=\id\) pour le quotient qui introduit une torsion ternaire finie dans
\(K_m\).
\end{proposition}

\begin{proof}
Une section homomorphe enverrait un élément d'ordre \(3^r\) de \(K_m\) sur un
élément du même ordre dans la composante \(\mathbb Z_3\). Mais le groupe
additif \(\mathbb Z_3\) est sans torsion : si \(3^rx=0\), alors \(x=0\).
L'image ne peut pas se projeter de nouveau sur un élément non nul d'ordre
\(3^r\).
\end{proof}

L'obstruction ne dit pas « on ne peut pas choisir ». Elle dit : on ne peut pas choisir
en préservant simultanément la projection et la loi de groupe. L'information de la
retenue qui disparaît dans le quotient réapparaît comme obstruction à la section.

\begin{remark}[Séparation des résultats et portée]
Le continu profini et l'obstruction de Pell sont des
\emph{résultats établis dans les sections précédentes}. Les deux théorèmes d'ensemble des parties et le tableau des
choix sont des \emph{constructions explicites du présent développement}. Le falsificateur dense empêche
de transformer la recherche incomplète d'un sous-ensemble en un ensemble des parties déjà
reconstruit.
\end{remark}
